\sectionheading{5.39 Python \textenglish{\ttfamily for} Loop}

কোনো Sequence বা Iterable-এর প্রতিটি Item নিয়ে একবার করে কাজ করতে
অথবা নির্দিষ্ট সংখ্যকবার কোনো Code Block Execute করতে Python-এ
\textenglish{\ttfamily for} Loop ব্যবহার করা হয়।

%================================================
% Definition
%================================================

\booksubheading{\textenglish{\ttfamily for} Loop কী?}

\begin{definitionbox}

কোনো Iterable-এর Itemগুলো পর্যায়ক্রমে গ্রহণ করে প্রতিটি Item-এর জন্য
নির্দিষ্ট Code Block Execute করার Loop-কে
\textbf{\textenglish{\ttfamily for} Loop} বলে।

\end{definitionbox}

\begin{center}

\textbf{Iterable / Sequence}
\(\longrightarrow\)
\textbf{একটি Item গ্রহণ}
\(\longrightarrow\)
\textbf{Loop Body Execute}
\(\longrightarrow\)
\textbf{পরবর্তী Item}

\end{center}

উদাহরণ—

\begin{codebox}
\begin{english}
\ttfamily
for i in range(5):\\
\hspace*{4mm}print(i)
\end{english}
\end{codebox}

\textbf{Output:}

\begin{center}
\textenglish{\ttfamily 0 \quad 1 \quad 2 \quad 3 \quad 4}
\end{center}

এখানে
\textenglish{\ttfamily i}
হলো \textbf{Loop Variable}।

%================================================
% Syntax
%================================================

\booksubheading{\textenglish{\ttfamily for} Loop-এর Syntax}

\begin{codebox}
\begin{english}
\ttfamily
for variable in iterable:\\
\hspace*{4mm}statement
\end{english}
\end{codebox}

উদাহরণ—

\begin{codebox}
\begin{english}
\ttfamily
for x in [10, 20, 30]:\\
\hspace*{4mm}print(x)
\end{english}
\end{codebox}

\textbf{Output:}

\begin{center}
\textenglish{\ttfamily 10 \quad 20 \quad 30}
\end{center}

\begin{tabularx}
{\linewidth}
{
|>{\bfseries\RaggedRight\arraybackslash}p{0.27\linewidth}
|Y|
}

\hline
\rowcolor{softBlue}

\multicolumn{1}{|c|}{\textbf{অংশ}}
&
\multicolumn{1}{c|}{\textbf{কাজ}}
\\
\hline

\textenglish{\ttfamily for}
&
Loop শুরু করার Keyword।
\\
\hline

\textenglish{\ttfamily x}
&
Loop Variable।
\\
\hline

\textenglish{\ttfamily in}
&
Iterable থেকে একে একে Item গ্রহণে ব্যবহৃত Keyword।
\\
\hline

\textenglish{\ttfamily [10,20,30]}
&
Iterable।
\\
\hline

\textenglish{\ttfamily :}
&
Loop Header-এর শেষ নির্দেশ করে।
\\
\hline

\textenglish{\ttfamily print(x)}
&
Loop Body।
\\
\hline

\end{tabularx}

%================================================
% Iterable
%================================================

\booksubheading{Iterable কী?}

যে Object-এর Itemগুলো একে একে গ্রহণ করে Loop চালানো যায় তাকে
\textbf{Iterable} বলা হয়।

Python-এ সাধারণ Iterable—

\begin{center}

\textbf{String}
\qquad
\textbf{List}
\qquad
\textbf{Tuple}
\qquad
\textbf{Set}
\qquad
\textbf{Dictionary}
\qquad
\textenglish{\ttfamily range()}

\end{center}

\begin{examplebox}{String-এর উপর Loop}

\begin{codebox}
\begin{english}
\ttfamily
for ch in "ICT":\\
\hspace*{4mm}print(ch)
\end{english}
\end{codebox}

\textbf{Output:}

\begin{center}
\textenglish{\ttfamily I \qquad C \qquad T}
\end{center}

\end{examplebox}

%================================================
% Loop Variable
%================================================

\booksubheading{Loop Variable কীভাবে কাজ করে?}

\begin{codebox}
\begin{english}
\ttfamily
for i in [10, 20, 30]:\\
\hspace*{4mm}print(i)
\end{english}
\end{codebox}

\begin{tabularx}
{\linewidth}
{
|>{\bfseries\centering\arraybackslash}p{0.30\linewidth}
|>{\centering\arraybackslash}Y|
}

\hline
\rowcolor{softBlue}

\textbf{Iteration}
&
\textbf{\textenglish{\ttfamily i}-এর Value}
\\
\hline

1
&
10
\\
\hline

2
&
20
\\
\hline

3
&
30
\\
\hline

\end{tabularx}

\begin{rememberbox}

\textenglish{\ttfamily for}
Loop Iterable থেকে নিজেই পরবর্তী Item নিয়ে Loop Variable-এ দেয়।

সাধারণত
\textenglish{\ttfamily while}
-এর মতো আলাদাভাবে
\textenglish{\ttfamily i += 1}
লিখে Loop Variable Update করতে হয় না।

\end{rememberbox}

%================================================
% range
%================================================

\booksubheading{\textenglish{\ttfamily range()} Function}

নির্দিষ্ট সংখ্যার ধারায়
\textenglish{\ttfamily for}
Loop চালাতে
\textenglish{\ttfamily range()}
খুব বেশি ব্যবহৃত হয়।

\begin{center}

\textenglish{\ttfamily range(stop)}

\qquad

\textenglish{\ttfamily range(start, stop)}

\qquad

\textenglish{\ttfamily range(start, stop, step)}

\end{center}

%================================================
% range stop + start stop
%================================================

\noindent
\begin{minipage}[t]{0.49\linewidth}

\begin{sideexamplebox}{১. \textenglish{\ttfamily range(stop)}}

\begin{sidecodebox}
\begin{english}
\ttfamily
for i in range(5):\\
\hspace*{4mm}print(i)
\end{english}
\end{sidecodebox}

\textbf{Output:}

\begin{center}
\textenglish{\ttfamily 0 \quad 1 \quad 2 \quad 3 \quad 4}
\end{center}

\textenglish{\ttfamily range(5)}
শুরু হয় 0 থেকে এবং 5 অন্তর্ভুক্ত হয় না।

\end{sideexamplebox}

\end{minipage}
\hfill
\begin{minipage}[t]{0.49\linewidth}

\begin{sideexamplebox}{২. \textenglish{\ttfamily range(start, stop)}}

\begin{sidecodebox}
\begin{english}
\ttfamily
for i in range(2, 6):\\
\hspace*{4mm}print(i)
\end{english}
\end{sidecodebox}

\textbf{Output:}

\begin{center}
\textenglish{\ttfamily 2 \quad 3 \quad 4 \quad 5}
\end{center}

Start অন্তর্ভুক্ত হয়, কিন্তু Stop অন্তর্ভুক্ত হয় না।

\end{sideexamplebox}

\end{minipage}

\begin{rememberbox}

\textbf{\color{deepOrange}MCQ Trap}

\begin{center}

\textenglish{\ttfamily range(2, 6)}
\(\longrightarrow\)
\textenglish{\ttfamily 2, 3, 4, 5}

\end{center}

\textbf{Start Included, Stop Excluded}

\end{rememberbox}

%================================================
% range step
%================================================

\booksubheading{৩. \textenglish{\ttfamily range(start, stop, step)}}

Step দ্বারা প্রতিবার Value কত করে পরিবর্তিত হবে তা নির্ধারিত হয়।

\begin{codebox}
\begin{english}
\ttfamily
for i in range(2, 11, 2):\\
\hspace*{4mm}print(i)
\end{english}
\end{codebox}

\textbf{Output:}

\begin{center}
\textenglish{\ttfamily 2 \quad 4 \quad 6 \quad 8 \quad 10}
\end{center}

এখানে—

\begin{center}

\textbf{Start}
\(\longrightarrow\)
\textenglish{\ttfamily 2}

\qquad

\textbf{Stop}
\(\longrightarrow\)
\textenglish{\ttfamily 11}

\qquad

\textbf{Step}
\(\longrightarrow\)
\textenglish{\ttfamily 2}

\end{center}

%================================================
% Default values
%================================================

\booksubheading{\textenglish{\ttfamily range()}-এর Default Value}

\begin{tabularx}
{\linewidth}
{
|>{\bfseries\RaggedRight\arraybackslash}p{0.35\linewidth}
|>{\centering\arraybackslash}Y|
}

\hline
\rowcolor{softBlue}

\textbf{অংশ}
&
\textbf{Default Value}
\\
\hline

Start
&
\textenglish{\ttfamily 0}
\\
\hline

Step
&
\textenglish{\ttfamily 1}
\\
\hline

Stop
&
অবশ্যই দিতে হয়
\\
\hline

\end{tabularx}

অর্থাৎ—

\begin{center}

\textenglish{\ttfamily range(5)}
\(\longrightarrow\)
0 থেকে শুরু
\(\longrightarrow\)
1 করে বৃদ্ধি

\end{center}

%================================================
% Forward + Reverse
%================================================

\booksubheading{Forward ও Reverse Loop}

\noindent
\begin{minipage}[t]{0.49\linewidth}

\begin{sideexamplebox}{Forward Loop}

Positive Step ব্যবহার করলে Value সাধারণত সামনে দিকে বাড়ে।

\begin{sidecodebox}
\begin{english}
\ttfamily
for i in range(1, 6):\\
\hspace*{4mm}print(i)
\end{english}
\end{sidecodebox}

\textbf{Output:}

\begin{center}
\textenglish{\ttfamily 1 \quad 2 \quad 3 \quad 4 \quad 5}
\end{center}

\end{sideexamplebox}

\end{minipage}
\hfill
\begin{minipage}[t]{0.49\linewidth}

\begin{sideexamplebox}{Reverse Loop}

পেছনের দিকে যেতে Negative Step ব্যবহার করা হয়।

\begin{sidecodebox}
\begin{english}
\ttfamily
for i in range(5, 0, -1):\\
\hspace*{4mm}print(i)
\end{english}
\end{sidecodebox}

\textbf{Output:}

\begin{center}
\textenglish{\ttfamily 5 \quad 4 \quad 3 \quad 2 \quad 1}
\end{center}

\end{sideexamplebox}

\end{minipage}

\begin{rememberbox}

\textbf{\color{deepOrange}Reverse Loop}

\begin{center}

\textbf{Start বড়}
\(\longrightarrow\)
\textbf{Stop ছোট}
\(\longrightarrow\)
\textbf{Step Negative}

\end{center}

উদাহরণ:
\textenglish{\ttfamily range(5, 0, -1)}

\end{rememberbox}

%================================================
% Step zero
%================================================

\booksubheading{Step কি \textenglish{\ttfamily 0} হতে পারে?}

না। Step
\textenglish{\ttfamily 0}
দেওয়া যায় না।

\begin{codebox}
\begin{english}
\ttfamily
range(1, 10, 0)
\end{english}
\end{codebox}

এটি Error সৃষ্টি করবে।

\begin{rememberbox}

\textbf{\color{deepOrange}MCQ তথ্য:}
\quad
\textenglish{\ttfamily range()}
-এর Step কখনো
\textenglish{\ttfamily 0}
হতে পারে না।

\end{rememberbox}

%================================================
% Range summary table
%================================================

\booksubheading{\textenglish{\ttfamily range()} এক নজরে}

\begin{tabularx}
{\linewidth}
{
|>{\bfseries\RaggedRight\arraybackslash}p{0.38\linewidth}
|Y|
}

\hline
\rowcolor{softBlue}

\textbf{Expression}
&
\textbf{Values}
\\
\hline

\textenglish{\ttfamily range(5)}
&
\textenglish{\ttfamily 0, 1, 2, 3, 4}
\\
\hline

\textenglish{\ttfamily range(1, 5)}
&
\textenglish{\ttfamily 1, 2, 3, 4}
\\
\hline

\textenglish{\ttfamily range(1, 6)}
&
\textenglish{\ttfamily 1, 2, 3, 4, 5}
\\
\hline

\textenglish{\ttfamily range(2, 11, 2)}
&
\textenglish{\ttfamily 2, 4, 6, 8, 10}
\\
\hline

\textenglish{\ttfamily range(5, 0, -1)}
&
\textenglish{\ttfamily 5, 4, 3, 2, 1}
\\
\hline

\textenglish{\ttfamily range(10, 0, -2)}
&
\textenglish{\ttfamily 10, 8, 6, 4, 2}
\\
\hline

\end{tabularx}

%================================================
% String List Tuple
%================================================

\booksubheading{String, List ও Tuple-এর উপর \textenglish{\ttfamily for} Loop}

\noindent
\begin{minipage}[t]{0.32\linewidth}

\begin{sideexamplebox}{String}

\begin{sidecodebox}
\begin{english}
\ttfamily
for ch in "ICT":\\
\hspace*{4mm}print(ch)
\end{english}
\end{sidecodebox}

Output:

\begin{center}
\textenglish{\ttfamily I \quad C \quad T}
\end{center}

\end{sideexamplebox}

\end{minipage}
\hfill
\begin{minipage}[t]{0.32\linewidth}

\begin{sideexamplebox}{List}

\begin{sidecodebox}
\begin{english}
\ttfamily
for x in [10, 20, 30]:\\
\hspace*{4mm}print(x)
\end{english}
\end{sidecodebox}

Output:

\begin{center}
\textenglish{\ttfamily 10 \quad 20 \quad 30}
\end{center}

\end{sideexamplebox}

\end{minipage}
\hfill
\begin{minipage}[t]{0.32\linewidth}

\begin{sideexamplebox}{Tuple}

\begin{sidecodebox}
\begin{english}
\ttfamily
for x in (2, 4, 6):\\
\hspace*{4mm}print(x)
\end{english}
\end{sidecodebox}

Output:

\begin{center}
\textenglish{\ttfamily 2 \quad 4 \quad 6}
\end{center}

\end{sideexamplebox}

\end{minipage}

%================================================
% Dictionary
%================================================

\booksubheading{Dictionary-এর উপর \textenglish{\ttfamily for} Loop}

Dictionary-এর উপর সরাসরি
\textenglish{\ttfamily for}
Loop চালালে সাধারণত তার \textbf{Key} পাওয়া যায়।

\begin{codebox}
\begin{english}
\ttfamily
student = \{"name": "Rafi", "age": 18\}\\
for x in student:\\
\hspace*{4mm}print(x)
\end{english}
\end{codebox}

\textbf{Output:}

\begin{center}
\textenglish{\ttfamily name \qquad age}
\end{center}

\begin{rememberbox}

\textbf{\color{deepOrange}MCQ তথ্য}

\begin{center}

\textbf{Dictionary-এর উপর সরাসরি}
\textenglish{\ttfamily for}
\textbf{Loop}
\(\longrightarrow\)
\textbf{Key-এর উপর Iteration}

\end{center}

\end{rememberbox}

%================================================
% 1 to 10 + even odd
%================================================

\booksubheading{সংখ্যার Range-এর উপর \textenglish{\ttfamily for} Loop}

\noindent
\begin{minipage}[t]{0.32\linewidth}

\begin{sideexamplebox}{১ থেকে ১০}

\begin{sidecodebox}
\begin{english}
\ttfamily
for i in range(1, 11):\\
\hspace*{4mm}print(i)
\end{english}
\end{sidecodebox}

Stop
\textenglish{\ttfamily 11}
অন্তর্ভুক্ত নয়।

\end{sideexamplebox}

\end{minipage}
\hfill
\begin{minipage}[t]{0.32\linewidth}

\begin{sideexamplebox}{Even}

\begin{sidecodebox}
\begin{english}
\ttfamily
for i in range(2, 11, 2):\\
\hspace*{4mm}print(i)
\end{english}
\end{sidecodebox}

Output:

\begin{center}
\textenglish{\ttfamily 2,4,6,8,10}
\end{center}

\end{sideexamplebox}

\end{minipage}
\hfill
\begin{minipage}[t]{0.32\linewidth}

\begin{sideexamplebox}{Odd}

\begin{sidecodebox}
\begin{english}
\ttfamily
for i in range(1, 11, 2):\\
\hspace*{4mm}print(i)
\end{english}
\end{sidecodebox}

Output:

\begin{center}
\textenglish{\ttfamily 1,3,5,7,9}
\end{center}

\end{sideexamplebox}

\end{minipage}

%================================================
% Even using condition
%================================================

\begin{examplebox}{Condition ব্যবহার করেও Even Number বের করা যায়}

\begin{codebox}
\begin{english}
\ttfamily
for i in range(1, 11):\\
\hspace*{4mm}if i \% 2 == 0:\\
\hspace*{8mm}print(i)
\end{english}
\end{codebox}

\end{examplebox}

%================================================
% Sum
%================================================

\booksubheading{Loop ব্যবহার করে Sum নির্ণয়}

১ থেকে ৫ পর্যন্ত সংখ্যার যোগফল—

\begin{codebox}
\begin{english}
\ttfamily
total = 0\\[0.5mm]
for i in range(1, 6):\\
\hspace*{4mm}total += i\\[0.5mm]
print(total)
\end{english}
\end{codebox}

\textbf{Output:}
\textenglish{\ttfamily 15}

\begin{center}

\textenglish{\ttfamily 0}
\(\longrightarrow\)
\textenglish{\ttfamily +1}
\(\longrightarrow\)
\textenglish{\ttfamily +2}
\(\longrightarrow\)
\textenglish{\ttfamily +3}
\(\longrightarrow\)
\textenglish{\ttfamily +4}
\(\longrightarrow\)
\textenglish{\ttfamily +5}
\(\longrightarrow\)
\textenglish{\ttfamily 15}

\end{center}

\begin{tabularx}
{\linewidth}
{
|>{\bfseries\RaggedRight\arraybackslash}p{0.31\linewidth}
|Y|
}

\hline
\rowcolor{softBlue}

\textbf{অংশ}
&
\textbf{কাজ}
\\
\hline

\textenglish{\ttfamily total = 0}
&
যোগফলের প্রাথমিক Value।
\\
\hline

\textenglish{\ttfamily i}
&
Loop Variable।
\\
\hline

\textenglish{\ttfamily total += i}
&
প্রতিটি Iteration-এ Sum Update করে।
\\
\hline

\end{tabularx}

%================================================
% Iteration count
%================================================

\booksubheading{Loop কতবার চলবে?}

\noindent
\begin{minipage}[t]{0.49\linewidth}

\begin{sideexamplebox}{৫ Iteration}

\begin{sidecodebox}
\begin{english}
\ttfamily
for i in range(5):\\
\hspace*{4mm}print(i)
\end{english}
\end{sidecodebox}

Values—

\begin{center}
\textenglish{\ttfamily 0,1,2,3,4}
\end{center}

মোট ৫টি Iteration।

\end{sideexamplebox}

\end{minipage}
\hfill
\begin{minipage}[t]{0.49\linewidth}

\begin{sideexamplebox}{৪ Iteration}

\begin{sidecodebox}
\begin{english}
\ttfamily
for i in range(1, 5):\\
\hspace*{4mm}print(i)
\end{english}
\end{sidecodebox}

Values—

\begin{center}
\textenglish{\ttfamily 1,2,3,4}
\end{center}

মোট ৪টি Iteration।

\end{sideexamplebox}

\end{minipage}

\begin{rememberbox}

\textbf{\color{deepOrange}ভুল করো না}

\textenglish{\ttfamily range()}
-এর Stop Value গণনায় অন্তর্ভুক্ত হয় না।

\begin{center}

\textenglish{\ttfamily range(1,5)}
\(\longrightarrow\)
\textenglish{\ttfamily 1,2,3,4}

\end{center}

\end{rememberbox}

%================================================
% Indentation
%================================================

\booksubheading{Loop Body ও Indentation}

Indentation দিয়েই বোঝা যায় কোন Statement Loop Body-এর অংশ।

\noindent
\begin{minipage}[t]{0.49\linewidth}

\begin{sideexamplebox}{দুটি Statement Loop Body-তে}

\begin{sidecodebox}
\begin{english}
\ttfamily
for i in range(3):\\
\hspace*{4mm}print(i)\\
\hspace*{4mm}print("ICT")
\end{english}
\end{sidecodebox}

দুটি
\textenglish{\ttfamily print()}
-ই প্রতিটি Iteration-এ Execute হবে।

\end{sideexamplebox}

\end{minipage}
\hfill
\begin{minipage}[t]{0.49\linewidth}

\begin{sideexamplebox}{একটি Statement Loop-এর বাইরে}

\begin{sidecodebox}
\begin{english}
\ttfamily
for i in range(3):\\
\hspace*{4mm}print(i)\\[0.5mm]
print("ICT")
\end{english}
\end{sidecodebox}

\textenglish{\ttfamily print("ICT")}
মাত্র একবার Execute হবে।

\end{sideexamplebox}

\end{minipage}

\begin{rememberbox}

\textbf{\color{deepOrange}MCQ Trap:}
\quad
Python-এ
\textbf{Indentation}
দিয়ে Loop Body নির্ধারিত হয়।

\end{rememberbox}

%================================================
% Nested for
%================================================

\booksubheading{Nested \textenglish{\ttfamily for} Loop-এর প্রাথমিক ধারণা}

একটি
\textenglish{\ttfamily for}
Loop-এর ভিতরে আরেকটি
\textenglish{\ttfamily for}
Loop থাকলে তাকে
\textbf{Nested for Loop} বলে।

\begin{examplebox}{উদাহরণ}

\begin{codebox}
\begin{english}
\ttfamily
for i in range(2):\\
\hspace*{4mm}for j in range(3):\\
\hspace*{8mm}print(i, j)
\end{english}
\end{codebox}

এখানে—

\begin{center}

\textbf{Outer Loop}
\(\longrightarrow\)
2 বার

\qquad

\textbf{Inner Loop}
\(\longrightarrow\)
প্রতিবার 3 বার

\end{center}

তাই
\textenglish{\ttfamily print()}
মোট—

\begin{center}

\(2 \times 3 = 6\) বার Execute হবে।

\end{center}

\end{examplebox}

\begin{rememberbox}

Nested Loop পরবর্তী
\textbf{Loop Pattern}
Topic-এ বিশেষভাবে ব্যবহৃত হবে।

\end{rememberbox}

%================================================
% for else
%================================================

\booksubheading{\textenglish{\ttfamily for ... else}}

Python-এ
\textenglish{\ttfamily for}
Loop-এর পরে Optional
\textenglish{\ttfamily else}
ব্যবহার করা যায়।

\begin{codebox}
\begin{english}
\ttfamily
for i in range(3):\\
\hspace*{4mm}print(i)\\
else:\\
\hspace*{4mm}print("Loop Finished")
\end{english}
\end{codebox}

\textbf{Output:}

\begin{codebox}
\begin{english}
\ttfamily
0\\
1\\
2\\
Loop Finished
\end{english}
\end{codebox}

Loop স্বাভাবিকভাবে শেষ হলে
\textenglish{\ttfamily else}
Block Execute হয়।

\begin{rememberbox}

\textenglish{\ttfamily break}
দিয়ে Loop বন্ধ হলে
\textenglish{\ttfamily for ... else}
-এর আচরণ ভিন্ন হয়।

এটি
\textbf{5.41 Loop Control Statement}
-এ বিস্তারিত আলোচনা হবে।

\end{rememberbox}

%================================================
% break continue preview
%================================================

\booksubheading{\textenglish{\ttfamily break} ও \textenglish{\ttfamily continue}-এর প্রাথমিক ধারণা}

\begin{tabularx}
{\linewidth}
{
|>{\bfseries\RaggedRight\arraybackslash}p{0.27\linewidth}
|Y|
}

\hline
\rowcolor{softBlue}

\textbf{Statement}
&
\textbf{প্রাথমিক কাজ}
\\
\hline

\textenglish{\ttfamily break}
&
Loop সঙ্গে সঙ্গে বন্ধ করে।
\\
\hline

\textenglish{\ttfamily continue}
&
বর্তমান Iteration-এর বাকি অংশ Skip করে পরবর্তী Iteration-এ যায়।
\\
\hline

\end{tabularx}

\begin{rememberbox}

\begin{center}

\textenglish{\ttfamily break}
\(\longrightarrow\)
\textbf{Loop বন্ধ}

\qquad

\textenglish{\ttfamily continue}
\(\longrightarrow\)
\textbf{Current Iteration Skip}

\end{center}

বিস্তারিত:
\textbf{5.41 Loop Control Statement}

\end{rememberbox}

%================================================
% When to use
%================================================

\booksubheading{\textenglish{\ttfamily for} Loop কখন ব্যবহার করা সুবিধাজনক?}

\begin{compacttwocolpoints}

\item নির্দিষ্ট Range-এর উপর Loop চালাতে।

\item String-এর Character Process করতে।

\item List, Tuple, Set ইত্যাদির Item Process করতে।

\item Collection-এর প্রতিটি Item নিয়ে কাজ করতে।

\item নির্দিষ্ট সংখ্যক Iteration প্রয়োজন হলে।

\item Sum, Count, Multiplication Table ইত্যাদি তৈরি করতে।

\end{compacttwocolpoints}

%================================================
% Key features
%================================================

\booksubheading{\textenglish{\ttfamily for} Loop-এর গুরুত্বপূর্ণ বৈশিষ্ট্য}

\begin{tabularx}
{\linewidth}
{
|>{\bfseries\RaggedRight\arraybackslash}p{0.34\linewidth}
|Y|
}

\hline
\rowcolor{softBlue}

\textbf{বিষয়}
&
\textbf{তথ্য}
\\
\hline

Keyword
&
\textenglish{\ttfamily for}
\\
\hline

Item গ্রহণ
&
\textenglish{\ttfamily in}
\\
\hline

Code Block
&
Indentation দ্বারা নির্ধারিত
\\
\hline

Numeric Range
&
\textenglish{\ttfamily range()}
\\
\hline

Default Start
&
\textenglish{\ttfamily 0}
\\
\hline

Default Step
&
\textenglish{\ttfamily 1}
\\
\hline

Stop
&
Excluded
\\
\hline

Step
&
Positive বা Negative হতে পারে
\\
\hline

Step
\textenglish{\ttfamily 0}
&
গ্রহণযোগ্য নয়
\\
\hline

\end{tabularx}

%================================================
% Common Mistakes
%================================================

\booksubheading{সাধারণ ভুল}

\begin{rememberbox}

\begin{itemize}

\item
\textenglish{\ttfamily for}
Header-এর শেষে
\textenglish{\ttfamily :}
না দেওয়া।

\item Loop Body-তে সঠিক Indentation না দেওয়া।

\item
\textenglish{\ttfamily range()}
-এর Stop Value Included মনে করা।

\item
\textenglish{\ttfamily range(1,5)}
-কে ১ থেকে ৫ পর্যন্ত মনে করা।

\item Reverse Loop-এ Negative Step না দেওয়া।

\item
\textenglish{\ttfamily range()}
-এ Step হিসেবে
\textenglish{\ttfamily 0}
দেওয়া।

\item Loop কতবার Execute হবে তা ভুল হিসাব করা।

\item Loop Variable-কে অপ্রয়োজনীয়ভাবে নিজে Update করা।

\item Loop-এর বাইরে থাকা Statement-কে Loop Body মনে করা।

\item
\textenglish{\ttfamily for}
ও
\textenglish{\ttfamily while}
-এর ব্যবহার গুলিয়ে ফেলা।

\end{itemize}

\end{rememberbox}

%================================================
% Final Reminder
%================================================

\begin{rememberbox}

\textbf{\color{deepOrange}মনে রাখবে}

\begin{tabularx}
{\linewidth}
{
>{\bfseries\RaggedRight\arraybackslash}p{0.25\linewidth}
Y
>{\bfseries\RaggedRight\arraybackslash}p{0.25\linewidth}
Y
}

\textenglish{\ttfamily for}
&
Iterable-এর Item নেয়
&
\textenglish{\ttfamily in}
&
Item গ্রহণের Keyword
\\[0.7mm]

\textenglish{\ttfamily range(stop)}
&
0 থেকে stop-1
&
Default Start
&
\textenglish{\ttfamily 0}
\\[0.7mm]

Default Step
&
\textenglish{\ttfamily 1}
&
Stop
&
Excluded
\\[0.7mm]

Positive Step
&
Forward
&
Negative Step
&
Reverse
\\[0.7mm]

Step
&
\textenglish{\ttfamily 0}
হতে পারে না
&
Dictionary
&
সরাসরি Loop-এ Key
\\[0.7mm]

Loop Body
&
Indentation দ্বারা
&
Nested Loop
&
Loop-এর ভিতরে Loop
\\

\end{tabularx}

\medskip

\begin{center}

\textenglish{\ttfamily range(5)}
\(\longrightarrow\)
\textenglish{\ttfamily 0,1,2,3,4}

\medskip

\textenglish{\ttfamily range(1,6)}
\(\longrightarrow\)
\textenglish{\ttfamily 1,2,3,4,5}

\medskip

\textenglish{\ttfamily range(2,11,2)}
\(\longrightarrow\)
\textenglish{\ttfamily 2,4,6,8,10}

\medskip

\textenglish{\ttfamily range(5,0,-1)}
\(\longrightarrow\)
\textenglish{\ttfamily 5,4,3,2,1}

\medskip

\textbf{Start Included}
\qquad
\textbf{Stop Excluded}

\end{center}

\end{rememberbox}

\small

%================================================
% 5.40 Python while Loop
%================================================

